% --- MAGIC COMMENTS: CONFIGURAZIONE EDITOR --- (vedi main.tex per riferimenti e spiegazioni) ---
% !TEX root = ../../main.tex

% ______________________________________________________ PACKAGES ______________________________________________________

% Funzionalità aggiuntive o modifiche del comportamento standard di LaTeX

% --- PACCHETTI relativi a FUNZIONI BASE ---
\usepackage{comment}                % [EXSTENTION] consente la creazione di commenti multilinea sfruttando l'enviroment
                                    % "comment" con il comando \begin. Nota che il contenuto all'internod di questi
                                    % blocchi di testo, non iniziando con "%", non è ignorato dal Parser di LTeX+, che
                                    % quindi effettuerà la correzione grammaticale e ortografica anche su questi blocchi

% --- PACCHETTI FORMATTAZIONE ---
% Nota: quelli "commentati" sono già aggiunti dalla custom class HKNdocument. Li lascio perchè molto importanti.

% \usepackage{lmodern}              % [PLUG & PLAY] [OVERRIDE] font vettoriali 100%, evita font accentati bitmap
                                    % sgranati. Pacchetto plug & play non c'è altro da usare.
% \usepackage[T1]{fontenc}          % [PLUG & PLAY] [OVERRIDE] aggiunge lettere accentate reali: LaTeX scriverà le
                                    % lettere accentate come oggetti singoli, e non come caratteri modificati.
                                    % Fondamentale per copia e incolla dal pdf compilato.
% \usepackage[utf8]{inputenc}       % [PLUG & PLAY] [OVERRIDE] consente a LaTeX di leggere le lettere accentate,
                                    % senza dover scrivere ad esempio "\'e" al posto di "è". Di solito è caricato
                                    % nativamente nelle versioni recenti di LaTeX, ma meglio specificare.
% \usepackage[italian]{babel}       % [PLUG & PLAY] [OVERRIDE] traduzione strutturale (termini associati al documento:
                                    % capitolo, sezione, eccetera) in italiano, e sillabazione italiana.

% --- PACCHETTI PER SCRITTURA DI MACRO AVANZATE ---
\usepackage{xparse}                 % [EXTENSION] Pacchetto che fornisce comandi avanzati per la definizione di macro,
                                    % come \NewDocumentCommand, grazie alla possibilità di definire argomenti delle
                                    % macro tramite "signature" flessibili.
                                    % Utilizzato nella macro \uchapter definita in 02_CustomMacros.tex.
                                    % NOTA SU XPARSE E LATEX3: Tecnicamente NON è necessario aggiungere:
                                    % - \RequirePackage{xparse} all'inizio della classe, oppure
                                    % - \usepackage{xparse} nel preambolo del main.tex
                                    % Questo perché dal 2020, xparse e le funzioni di LaTeX3 sono state integrate
                                    % nativamente nel kernel standard di LaTeX2e. Tuttavia, è good practice includerlo
                                    % esplicitamente, per garantire la retrocompatibilità con vecchie distribuzioni
                                    % TeX antecedenti al 2020. In ogni caso non causa problemi.
                                    % Comunque, proprio per questo stesso motivo di integrazione nativa, ad esempio non
                                    % bisogna modificare l'istruzione \NeedsTeXFormat{LaTeX2e} (es. in "LaTeX3")
                                    % presente nella custom class, poiché LaTeX3 non è mai stato rilasciato come
                                    % formato a sé stante, ma è un ambiente di programmazione integrato dentro LaTeX2e.

% --- PACCHETTI PER BIBLIOGRAFIA E PER GESTIONE SMART DEGLI APICI ---
\usepackage[backend=biber, autolang=other, style=ieee]{biblatex}
\begin{comment} BibLaTeX - PACCHETTO DI GESTIONE AVANZATA DELLA BIBLIOGRAFIA

    BibLaTeX è un pacchetto per la gestione avanzata della bibliografia,
    standard moderno per la bibliografia in LaTeX.
    - backend=biber: Specifica il motore bibliografico da utilizzare per
      l'elaborazione dei dati contenuti nel database bibliografico (.bib),
      in questo caso generato con Zotero.
      Nelle distribuzioni moderne di TeX, in precisione TeX Live, è già
      impostato di default, tuttavia esplicitarlo è la best practice, e
      garantisce la compatibilità anche con distribuzioni più datate.
      Offre pieno supporto nativo al formato UTF-8 e algoritmi di ordinamento e
      parsing molto più avanzati del legacy BibTeX.

    - autolang=other: Risolve i problemi di sillabazione della lingua nella
      bibliografia, nel caso le fonti citate abbiano più lingue diverse.
      È fondamentale che il campo "langid" nel database .bib sia esplicitato
      correttamente, in quanto è il campo usato per selezionare come viene
      trattata ogni voce della bibliografia.
      Zotero trasforma la lingua in langid automaticamente, se è scritta
      correttamente (codici ISO, o per esteso in inglese).
      In precisione "other" crea un ambiente linguistico isolato per la singola
      riga della bibliografia: adatta sia le regole di sillabazione che di
      formattazione e di terminologia funzione della localizzazione (esempio:
      Edition vs Edizione, a cura di vs edited by, eccetera) garantendo le
      migliori impostazioni tipografiche per ogni voce.
      Come alternativa si può usare "autolang=hyphen", che si limita solo alla
      sillabazione multilingua, senza adattare il resto.

    - style=ieee: applica lo standard bibliografico dell'Institute of Electrical
      and Electronics Engineers (numerico con parentesi quadre, titoli
      formattati e abbreviazioni specifiche). In precisione:
        * Citazioni testuali: numeriche tra parentesi quadre (es. [1]),
        con accorpamento automatico per fonti multiple (es. [1]-[3]).
        * Nomenclatura: autori indicati con iniziali del nome e cognome per
          esteso (es. J. C. Maxwell).
        * Tipografia dei titoli: titoli di articoli o paper racchiusi tra
          virgolette, mentre i nomi di riviste, manuali e libri sono formattati
          in corsivo.
        * Ordinamento bibliografico: rigorosamente sequenziale, basato
          sull'esatto ordine di comparsa nel testo (sorting 'none').
          Se si forza la stampa delle fonti non citate, vengono elencate dopo
          quelle citate, nell'ordine di comparsa nel database .bib.
        * Dettagli tecnici: standardizzazione automatica della punteggiatura
        e delle abbreviazioni specifiche (Vol., No., pp.).
      È la scelta coerente e per l'identità del Mu Nu Chapter HKN.
      (Nota: richiede che nel sistema sia installato 'biblatex-ieee', incluso
      con TexLive. Se si usa una distribuzione diversa, lo si può installare da
      terminale col comando "tlmgr install biblatex-ieee").

\end{comment}

\usepackage{csquotes}               % Richiesto da biblatex per la gestione intelligente delle virgolette e degli apici
                                    % nelle citazioni bibliografiche. Garantisce che le citazioni siano formattate
                                    % correttamente in base alla lingua del documento (es. virgolette italiane « »).
                                    % Fornisce inoltre utilissimi comandi per la gestione degli apici, come:
                                    % - \enquote{text} per riportare il discorso diretto in funzione della lingua
                                    %   selezionata con babel (es: sfrutta le caporali in italiano « »), ed in funzinone
                                    %   del contesto semantico (es. se la citazione è all'interno di un'altra citazione,
                                    %   la citazione di secondo livello, nidificata con le virgolette doppie normali "".
                                    % - \enquote*{text} ignora il primo livello è usa sempre le virgolette doppie
                                    %   normali "".

\addbibresource{chapters/00.2_Bibliography/HKN_ElnBiliography.bib}
% File (Database bibliografico) contentnente tutte le fonti bibliografiche di interesse per il progetto.
% ATTENZIONE: il file viene importato nel main.tex con il comando \input. Tieni a mente che TeX è un linguaggio
% Espansore di macro: quindi il codice è come se fosse scritto nel main.tex, e per questo motivo è importante che i
% percorsi siano riferti alla posizione del main.tex.


% --- PACCHETTI PER FISICA ---
\usepackage{siunitx}                % [EXTENSION] Per le unità di misura (es. \si{\ohm})

% --- PACCHETTI PER GRAFICA ---
% \usepackage{circuitikz}             % [EXTENSION] Per disegnare i circuiti elettrici
% \usepackage{pgfplots}               % [EXTENSION] Per i grafici matematici e di Bode

\begin{comment}     % --- RIEPILOGO PACCHETTI INCLUSI NELLA CLASSE "HKNdocument" ---

    RIEPILOGO COMPLETO: PACCHETTI CARICATI DALLA CLASSE HKNdocument E _preamble.tex
    NON inserirli di nuovo in main.tex con \usepackage per evitare conflitti!

    --- TIPOGRAFIA E LINGUA BASE (da HKNdocument.cls) ---
    \usepackage[T1]{fontenc}                    % [PLUG & PLAY] [OVERRIDE]  Imposta la codifica europea dei font,
                                                % fondamentale per stampare correttamente le lettere accentate italiane
                                                % nel PDF e permettere il copia-incolla.
    \usepackage[utf8]{inputenc}                 % [PLUG & PLAY] [OVERRIDE]  Permette di digitare direttamente da
                                                % tastiera caratteri speciali e accenti (es. à, è, ì) riconoscendo la
                                                % codifica UTF-8 (unicode) direttamente dal file sorgente.
                                                % Utile per non impazzire con i comandi per aggiungere apici.
    \usepackage{lmodern}                        % [PLUG & PLAY] [OVERRIDE]  Carica i font vettoriali "Latin Modern",
                                                % essenziali per mantenere il testo nitido e non sgranato anche ad alti
                                                % livelli di zoom.
    \usepackage[\DocumentLanguage]{babel}       % [PLUG & PLAY] [OVERRIDE]  Gestisce le regole tipografiche, le
                                                % traduzioni automatiche (es. "Capitolo") e la corretta sillabazione per
                                                % la lingua scelta (italiano o inglese).
                                                % NOTA: il "\DocumentLanguage" è l'opzione passata alla document class
                                                % custom HKNdocument, all'inizio del main.tex

    --- LAYOUT E IMPAGINAZIONE (da HKNdocument.cls) ---
    \usepackage{geometry}           % [PLUG & PLAY] [OVERRIDE]  Configura fisicamente la pagina, impostando il formato
                                    % A4 e forzando margini uniformi di 3 cm su tutti i lati per una lettura ariosa.
    \usepackage{parskip}            % [PLUG & PLAY] [OVERRIDE]  Rimuove l'indentazione o rientro della prima riga
                                    % (considerata antiestetica), e aggiunge automaticamente spazio verticale tra i
                                    % paragrafi.
    \usepackage{fancyhdr}           % [EXTENSION]   [OVERRIDE]  Permette la creazione di intestazioni e piè di pagina
                                    % personalizzati (usati dalla classe per inserire loghi, pagina e titolo del
                                    % capitolo).

    --- MATEMATICA E FISICA (da _preamble.tex) ---
    \usepackage{amsmath}            % [EXTENSION] Il pacchetto fondamentale dell'American Mathematical Society per
                                    % scrivere equazioni complesse, matrici e sistemi matematici.
    \usepackage{amssymb}            % [EXTENSION] Aggiunge un'enorme libreria di simboli matematici avanzati (es. i
                                    % simboli degli insiemi dei numeri reali o complessi).
    \usepackage{amsthm}             % [EXTENSION] Fornisce gli ambienti standard per definire teoremi, dimostrazioni e
                                    % lemmi matematici.
    \usepackage{mathtools}          % [EXTENSION] [EXTENSION] Un'estensione di amsmath che risolve alcuni bug grafici e
                                    % aggiunge strumenti avanzati per allineare le equazioni.
    \usepackage{mleftright}         % [PLUG & PLAY] [OVERRIDE] [EXTENSION]  Ottimizza e corregge la spaziatura
                                    % orizzontale automatica attorno alle parentesi che racchiudono frazioni o matrici
                                    % alte.

    --- TABELLE ED ELENCHI (da _preamble.tex) ---
    \usepackage{longtable}          % [EXTENSION] Permette la creazione di tabelle lunghe che si spezzano
                                    % automaticamente e in modo pulito su più pagine.
    \usepackage{array}              % [EXTENSION] [OVERRIDE] Estende le funzionalità delle tabelle, permettendo di
                                    % definire larghezze precise per le colonne e allineamenti complessi.
    \usepackage{booktabs}           % [EXTENSION] Standard editoriale per tabelle professionali: rimuove le linee
                                    % verticali e introduce linee orizzontali di diverso spessore (\toprule, \midrule,
                                    % \bottomrule).
    \usepackage{enumitem}           % [EXTENSION] [OVERRIDE] Offre il controllo totale sugli elenchi puntati e numerati,
                                    % permettendo di regolarne i margini e la spaziatura verticale.

    --- GRAFICA E COLORI (da HKNdocument.cls e _preamble.tex) ---
    \usepackage{xcolor}             % [EXTENSION] Il motore principale per definire, miscelare e applicare colori
                                    % personalizzati al testo e agli sfondi.
    \usepackage{tcolorbox}          % [EXTENSION] Un potente pacchetto grafico usato dalla classe per creare i riquadri
                                    % arrotondati e colorati di definizioni ed esercizi.
    \usepackage{graphicx}           % [EXTENSION] Il pacchetto standard e indispensabile per importare e ridimensionare
                                    % immagini esterne (JPG, PNG, PDF).
    \usepackage{svg}                % [EXTENSION] Permette di importare direttamente file vettoriali .svg (utilizzando
                                    % Inkscape in background per la conversione).
    \usepackage{float}              % [OVERRIDE] [EXTENSION] Aggiunge l'opzione [H] (exactly here) e [h] (here if
                                    % possible, otherwise as usual) per forzare tabelle e figure a rimanere nel punto
                                    % del testo in cui sono state scriaggiunte, senza lasciarle "galleggiare" (float)
                                    % alla deriva altrove.
    \usepackage{caption}            % [PLUG & PLAY] [OVERRIDE] [EXTENSION]  Usato per personalizzare i margini, i font e
                                    % l'allineamento delle didascalie sotto le immagini e le tabelle.
                                    % Come Fix auomatico "Plug & play" corregge uno storico problema di LaTeX relativo
                                    % alla spaziatura delle didascalie se impostate sopra le tabelle. Inoltre la custom
                                    % class hai impostato ad 1cm questo spazio.

    --- STILE E UTILITY (da HKNdocument.cls e _preamble.tex) ---
    \usepackage{hyperref}           % [PLUG & PLAY] [OVERRIDE] [EXTENSION] Rende il PDF interattivo: trasforma l'indice,
                                    % i riferimenti incrociati e gli URL in link cliccabili (impostati in blu o nero).
    \usepackage{listings}           % [EXTENSION] Permette di formattare ed evidenziare la sintassi di blocchi di codice
                                    % informatico.
                                    % NOTA: nella custom class è definito un TEMA CUSTOM, chiamato "hkn", per l'aspetto
                                    % del codice. Non è stato impostato come predefinito, va usato con
                                    % \begin{lstlisting}[style=hkn].
                                    % Si dodvrebbe correggere la classe per imposarlo come predefinito in generale,
                                    % basta aggiungere il comando \lstset{style=hkn}. Aggiunto nel file global settings.
    \usepackage{ulem}               % [EXTENSION] Fornisce comandi avanzati per sottolineare e barrare il testo.
    \usepackage{contour}            % [EXTENSION] Crea un contorno attorno ai caratteri, usato in combinazione con ulem
                                    % per creare sottolineature che non tagliano le lettere discendenti.

\end{comment}
